%!TEX root = ../../thesis.tex

\section{Overview}\label{sec:input-overview}

% NOTE: This chapter and \cref{chap:corpus-actions} use the SAME six sections in the
% SAME order, so the two formulations can be read as a diff. Keep them aligned when
% editing: Overview / Action Space / Observation / Reward and Episode Structure /
% Determinism / Relation to Prior Work.

% TODO: State the formulation in one paragraph: an action is one symbol of the input
% the program consumes --- a byte on stdin for the byte-fed targets, a protocol
% message for open62541 --- and an episode builds one test case left to right. Give
% the step loop as an algorithm, mirroring the layout of the corpus loop in
% \cref{sec:corpus-action-space} so the two algorithms can be compared line by line.
%
% Running example for this chapter: sequence (\cref{sec:target-sequence}), because
% episode-step position equals prefix depth, which makes every claim about credit
% assignment directly observable.
%
% Say up front what this formulation gives up: there is no corpus, no seed reuse and
% no retention decision --- everything the agent has learned lives in the network
% weights rather than in a population of inputs. That is the deliberate contrast with
% \cref{chap:corpus-actions}.
